\section*{अध्याय \ref{ch:b2:cell-thermodynamics} --- विद्युत-रासायनिक सेलों की ऊष्मागतिकी}

\begin{solution}{exo:b2:cell-thermodynamics:1}
\ce{H2 + 1/2O2 -> H2O(l)}, $n = 2$: $E^\circ = 237\,140/(2 \times 96\,485) =
\qty{1.229}{V}$।
\end{solution}

\begin{solution}{exo:b2:cell-thermodynamics:2}
$Q = 1.0 \times 3600 = \qty{3.6e3}{C}$; $n(\ce{Zn}) = 3600/(2 \times 96\,485) =
\qty{1.87e-2}{mol}$, अर्थात् \qty{1.22}{g}।
\end{solution}

\begin{solution}{exo:b2:cell-thermodynamics:3}
$\Delta_r G = -2 \times 96\,485 \times 1.10 = \qty{-212}{kJ/mol}$।
\qty{10.0}{g} के लिए, $10.0/65.38 = \qty{0.153}{mol}$: अधिक से अधिक \qty{32.5}{kJ}।
\end{solution}

\begin{solution}{exo:b2:cell-thermodynamics:4}
$U = (0.0592/2)\log(0.10/1.0 \times 10^{-3}) = \qty{0.059}{V}$; धनात्मक ध्रुव
(कैथोड) \qty{0.10}{mol/L} विलयन वाला इलेक्ट्रोड है।
\end{solution}

\begin{solution}{exo:b2:cell-thermodynamics:5}
$\Delta_r G^\circ = -2 \times 96\,485 \times 1.015 = \qty{-195.9}{kJ/mol}$;
$\Delta_r S^\circ = 2 \times 96\,485 \times (-4.0 \times 10^{-4}) =
\qty{-77.2}{J/(K.mol)}$; $\Delta_r H^\circ = -195.9 + 298 \times (-0.0772) =
\qty{-218.9}{kJ/mol}$।
\end{solution}

\begin{solution}{exo:b2:cell-thermodynamics:6}
(a) दिया गया कार्य \qty{195.9}{kJ}; ऊष्मा $T\Delta_r S = \qty{-23.0}{kJ}$: \qty{23.0}{kJ}
मुक्त। (b) कार्य $2 \times 96\,485 \times 0.90 = \qty{173.7}{kJ}$; मुक्त ऊष्मा
$-\Delta_r H - 173.7 = 218.9 - 173.7 = \qty{45.2}{kJ}$। अतिरिक्त
\qty{22.2}{kJ} धारा के कारण खोया गया कार्य है।
\end{solution}

\begin{solution}{exo:b2:cell-thermodynamics:7}
\ce{Fe^3+ + 3e- -> Fe}: $\Delta_r G^\circ = +\qty{4.7}{kJ/mol}$, $E^\circ =
-4700/(3 \times 96\,485) = \qty{-0.016}{V}$। $E_1^\circ = 0.769$ और
$E_2^\circ = \qty{-0.409}{V}$ के साथ: $(0.769 - 0.818)/3 = \qty{-0.016}{V}$।
\end{solution}

\begin{solution}{exo:b2:cell-thermodynamics:8}
$E^\circ = \qty{0.222}{V}$ (अध्याय की तरह)। $E = 0.222 - 0.0592\log 0.10 =
\qty{0.281}{V}$।
\end{solution}

\begin{solution}{exo:b2:cell-thermodynamics:9}
$\Delta_r G^\circ = 2(-318.30) = \qty{-636.6}{kJ}$, $n = 4$ के लिए: $E^\circ =
636\,600/(4 \times 96\,485) = \qty{1.65}{V}$। प्रति किलोग्राम ज़िंक, $1000/65.38 =
\qty{15.3}{mol}$, $15.3 \times 318.3 = \qty{4.87}{MJ}$, अर्थात् \qty{1.35}{kW.h}।
\end{solution}

\begin{solution}{exo:b2:cell-thermodynamics:10}
कार्बन $-2$ से $+4$ तक जाता है: $n = 6$। $\Delta_r H^\circ = -393.51 + 2(-285.83) +
239 = \qty{-726.2}{kJ/mol}$; $\Delta_r S^\circ = 213.8 + 2(69.95) - 127.2 - 1.5
\times 205.15 = \qty{-81.2}{J/(K.mol)}$; $\Delta_r G^\circ = -726.2 + 298.15 \times
0.0812 = \qty{-702.0}{kJ/mol}$। $E^\circ = 702\,000/(6 \times 96\,485) =
\qty{1.21}{V}$; अधिकतम दक्षता $702.0/726.2 = 0.967$। (वास्तविक मेथनॉल सेल
इस वोल्टता से बहुत नीचे काम करते हैं।)
\end{solution}

\begin{solution}{exo:b2:cell-thermodynamics:11}
$\Delta_r G = -2 \times 96\,485 \times 0.500 = \qty{-96.5}{kJ}$; $\Delta_r S = +\qty{57.9}{J/K}$;
$\Delta_r H = -96.5 + 298 \times 0.0579 = \qty{-79.2}{kJ}$। कार्य, \qty{96.5}{kJ},
$-\Delta_r H = \qty{79.2}{kJ}$ से अधिक है: सेल अपने परिवेश से $T\Delta_r S = \qty{17.3}{kJ}$
ऊष्मा अवशोषित करता है और उसे भी कार्य में बदल देता है। द्वितीय नियम
का उल्लंघन नहीं होता, क्योंकि अभिक्रिया की एन्ट्रॉपी उतनी ही बढ़ती है।
\end{solution}

\begin{solution}{exo:b2:cell-thermodynamics:12}
हर हाइड्रोजन इलेक्ट्रोड का $E = -0.0592\,\mathrm{pH}$ है: $U = 0.0592\,(\mathrm{pH}_s -
7.00) = 0.177$, $\mathrm{pH}_s = 10.0$। दोनों इलेक्ट्रोड एक ही युग्म के हैं:
मानक विभव कट जाता है, जैसा किसी भी \omterm{def:b2:cell-thermodynamics:concentration-cell}{सांद्रता सेल} में।
\end{solution}

\begin{solution}{pb:b2:cell-thermodynamics:1}
\textbf{1.} ऐनोड: \ce{H2 -> 2H+ + 2e-}; कैथोड: \ce{1/2O2 + 2H+ + 2e- -> H2O};
सेल: \ce{H2 + 1/2O2 -> H2O}।
\textbf{2.} दो।
\textbf{3.} $E^\circ = 237\,140/(2 \times 96\,485) = \qty{1.229}{V}$।
\textbf{4.} $228\,580/192\,970 = \qty{1.185}{V}$।
\textbf{5.} \qty{298}{K} पर $\Delta_r G^\circ$ में अंतर $\Delta_{\text{वाष्पन}}G^\circ(\ce{H2O}) = -228.58 + 237.14 =
\qty{8.56}{kJ/mol}$ है: वहाँ वाष्प द्रव से कम स्थायी है, और
वोल्टता $8560/192\,970 = \qty{0.044}{V}$ कम है।
\textbf{6.} $\Delta_r H^\circ = \qty{-285.83}{kJ/mol}$; $237.14/285.83 = 0.830$।
\textbf{7.} $\Delta_r S^\circ = 69.95 - 130.68 - 102.58 = \qty{-163.3}{J/(K.mol)}$।
\textbf{8.} $dE^\circ/dT = -163.3/192\,970 = \qty{-8.46e-4}{V/K}$।
\textbf{9.} $1.229 - 8.46 \times 10^{-4} \times 55 = \qty{1.182}{V}$।
\textbf{10.} \qty{400}{K} पर वही आकलन $1.229 - 8.46 \times 10^{-4}
\times 101.85 = \qty{1.143}{V}$ देता है; JANAF $221\,010/192\,970 = \qty{1.145}{V}$ देता है:
$\Delta_r S^\circ$ को स्थिर मानने की क़ीमत केवल \qty{2}{mV} है।
\textbf{11.} प्रति मोल $T\Delta_r S^\circ = 298.15 \times (-163.3) = \qty{-48.7}{kJ}$:
\qty{48.7}{kJ} मुक्त।
\textbf{12.} $\Delta_r G^\circ(353) \approx -237.14 + 55 \times 0.1633 =
\qty{-228.2}{kJ}$; $\Delta_r H^\circ$ के लगभग अपरिवर्तित रहने पर $228.2/285.8 = 0.80$।
\textbf{13.} $2 \times 96\,485 \times 0.70 = \qty{135.1}{kJ}$।
\textbf{14.} $135.1/285.83 = 0.47$; $135.1/237.14 = 0.57$।
\textbf{15.} $285.83 - 135.1 = \qty{150.7}{kJ}$।
\textbf{16.} \qty{48.7}{kJ} तो उत्क्रमणीय सेल भी मुक्त करता; शेष
\qty{102.0}{kJ} धारा के कारण हुई हानियों (प्रतिरोध और इलेक्ट्रोड
अधिविभव) से आते हैं।
\textbf{17.} एक सेल: $300/(2 \times 96\,485) = \qty{1.55e-3}{mol/s}$; पुंज:
$400 \times 1.55 \times 10^{-3} = \qty{0.622}{mol/s}$, \qty{1.25}{g/s}।
\textbf{18.} $400 \times 0.70 \times 300 = \qty{84}{kW}$।
\textbf{19.} $0.622 \times 150.7 = \qty{93.7}{kW}$: बिजली से अधिक ऊष्मा।
\textbf{20.} $5000/2.016 = \qty{2.48e3}{mol}$।
\textbf{21.} $2 \times 2480 \times 96\,485 = \qty{4.79e8}{C}$।
\textbf{22.} किसी सेल से गुज़रने वाला हर कूलॉम \qty{0.70}{J} देता है: $4.79 \times
10^8 \times 0.70 = \qty{3.35e8}{J}$।
\textbf{23.} $3.35 \times 10^8/3.6 \times 10^6 = \qty{93}{kW.h}$।
\textbf{24.} पूर्ण शक्ति पर $93/84 = \qty{1.1}{h}$।
\textbf{25.} प्रति सेल \qty{0.70}{V} पर \qty{5.0}{kg} हाइड्रोजन
$\boldsymbol{\approx \qty{93}{kW.h}}$ विद्युत ऊर्जा देती है।
\end{solution}
